Figure~\ref{fig:breakpoint} combines these runs with six random task samples per cell around the change (Appendix~\ref{app:transition}, Table~\ref{tab:transition}).
In these runs the change occurs between 30 and 35 locations (2.0--2.6 million arcs).
At 30 locations the direct root LP takes 570--1{,}300 seconds and the direct MIP usually finds a near-optimal schedule, but in up to two of the six samples per task count it still holds only its starting schedule at 1{,}800 seconds, and repeated runs of the same input can differ.
At 35 locations the root LP takes 1{,}140--1{,}740 seconds and the direct MIP holds only its starting schedule on 17 of 18 inputs, while the MIP over the CG routes gives the better schedule on all 18.
With 3{,}600 seconds for each MIP, on the rerun inputs (30 locations with 60--120 tasks, 35 locations with 70 and 120 tasks, and the 40--120-location inputs of Table~\ref{tab:locscaling}), the change moves by about five locations: at 35 locations the direct MIP then finds a schedule on 11 of 12 inputs (median gaps 0.04--0.42\%), whereas from 40 locations it still holds only its starting schedule on every input; the gaps of the MIP over the CG routes change by at most 0.9 percentage points.
Counting its CG phase, the MIP over the CG routes finished within 57 minutes on every input from 40 to 110 locations (67 minutes at 120), so on these inputs the direct MIP with 3{,}600 seconds had more total time than CG and still held only its starting schedule.
